\section{Hermitian forms and ordered fields}

Throughout this paragraph, K denotes a maximal ordered field (hence commutative and of characteristic zero; cf.~Chap.~VI, § 2), and we assume that we are in one of the following three cases:

1°) A = K, J is the identity;

2°) A is the field \(\mathrm{K}(i)\) , obtained by adjoining to K a square root i of -1, and, for every \(\lambda \in A\) , \(\overline{\lambda}\) is the conjugate of \(\lambda\) (chap.~II, § 7, no. 7).

3°) A is the quaternion algebra over K corresponding to the pair \((-1,-1)\) (or, as we shall say for short, the field of quaternions over K), and for every \(\lambda\in A\) , \(\overline{\lambda}\) is the conjugate of \(\lambda\) (cf.~chap.~II, § 7, no. 8 and chap.~VIII, § 11, no. 2).

If \(\Phi\) is a Hermitian form on E, so in all cases we have, \(\Phi(x, x) \in K\) for all \(x \in E\) , since \(\Phi(x, x) = \overline{\Phi(x, x)}\) .

